\documentclass[10pt,a4paper]{amsart}
\usepackage[margin=19mm]{geometry}
\usepackage{amsmath,amssymb,mathtools}
\usepackage[scr=rsfs,cal=euler]{mathalfa}
\usepackage{xfrac}
\usepackage[unicode,hidelinks,bookmarks=false]{hyperref}
\newcommand{\ie}{i.e.\ }
\newcommand{\cf}{cf.\ }
\newcommand{\resp}{respectively\ }
\newcommand{\Subsection}{\subsection}
\providecommand{\nobreakdash}{\nobreak\hbox{-}}
\setlength{\emergencystretch}{2em}
\multlinegap0pt
\newtheorem{theorem}{Theorem}
\newtheorem{proposition}{Proposition}
\def\R#1{\ensuremath{\mathbb{R}^{#1}}}
\def\Y{Z}
\DeclareMathOperator{\arccot}{arccot}
\def\tgamma{\widetilde{\gamma}}
\title{Concircular helices and concircular surfaces in Euclidean 3-space $\R3$}
\author{Pascual Lucas, Jos\'e Antonio Ortega-Yag\"ues}
\date{Source: arXiv:2601.19252v1 (CC BY 4.0)}
\begin{document}
\maketitle

\noindent\emph{Source and licence.} Concircular helices and concircular surfaces in Euclidean 3-space $\R3$. Version arXiv:2601.19252v1 (CC BY 4.0), distributed by arXiv under the Creative Commons Attribution 4.0 International licence, http://creativecommons.org/licenses/by/4.0/, as recorded in the arXiv licence metadata for this exact version and shown on the abstract page https://arxiv.org/abs/2601.19252v1. Full licence notice: \texttt{licenses/arxiv-2601.19252-CC-BY-4.0.txt}.\par\smallskip

\noindent\textit{Editorial scope.} Counted unit: the article's theorem th1 (source0.tex lines 498--500). The article's complete original proof of it is delivered with it (source0.tex lines 501--539, one \texttt{proof} environment of 39 lines). No mathematics is written, repaired or re-derived: the statement and the proof are the article's own lines, in the article's own notation, and no retained line is substituted or re-flowed.  Retained uncounted as the source-local context the proof needs: lines 142--144; lines 154--156; lines 166--172; lines 456--458; lines 483--486; lines 488--496; lines 490--494.  Nothing was omitted from any retained span, so the package carries no omission record.  No retained line contains a citation, so the wrapper carries no bibliography and no bibliography entry.  No retained line contains a figure environment, an image-inclusion command or drawing code, so no image is delivered and the package declares no asset dependency.  Editorial formatting only, in addition: an AMS \texttt{amsart} preamble that restates the article's own declarations and macros used by the retained lines, namely the environments \texttt{theorem}, \texttt{proposition} on separate counters, together with the standard packages those lines need.  The retained environments print in this document as Theorem 1, Proposition 1 in order of appearance, whereas the article prints the same items with the numbering of its own full document.  The complete native source of the article as distributed in the arXiv e-print for this exact version is bundled unchanged as \texttt{sources/193444/source0.tex}.
\begin{equation}\label{conc4}
\left(\frac{\rho'}{\kappa_\gamma(1+\rho^2)^{3/2}}\right)'= m\,\frac{\rho''}{\kappa_\gamma^2(1+\rho^2)^{5/2}},
\end{equation}

\begin{equation}\label{ecub}
b=\frac{1}{\rho'}\big(\mu+\lambda\kappa_\gamma(1+\rho^2)\big),
\end{equation}

\begin{theorem}\label{teoconc1}
Let $\gamma$ be an arclength parametrized curve, with $\kappa_\gamma>0$ and $\rho'\neq0$. $\gamma$ is a proper concircular helix if and only if its curvature $\kappa_\gamma$ and function $\rho=\tau_\gamma/\kappa_\gamma$ satisfy the following differential equation:
\[
\left(\frac{\rho'}{\kappa_\gamma(1+\rho^2)^{3/2}}\right)'= m\,\frac{\rho''}{\kappa_\gamma^2(1+\rho^2)^{5/2}},
\]
for a certain nonzero constant $m\in\R{}$. Moreover, a concircular vector field $Y$ for the proper concircular helix $\gamma$ is the extension of the vector field  $V=b\,D_\gamma+ \lambda N_\gamma$, where $b$ is the differentiable function given by (\ref{ecub}).
\end{theorem}

\begin{equation}\label{parY2}
\Y(u,v)=\alpha(u)+(v-u)\,T_\alpha(u),\quad v\neq u,
\end{equation}

\begin{align}
\kappa_\gamma(s) &= -\sin\omega(s)\,u'(s)\,\tau_\alpha(u(s))= \sin\omega(s)\,\omega'(s)\,(\tau_\alpha/\kappa_\alpha)(u(s)),\label{kg2}\\
\tau_\gamma(s) &= \cos\omega(s)\,u'(s)\,\tau_\alpha(u(s))= -\cos\omega(s)\,\omega'(s)\,(\tau_\alpha/\kappa_\alpha)(u(s)).\label{tg2}
\end{align}

\begin{proposition}\label{p3}
Let $M$ be a proper concircular surface parametrized by (\ref{parY2}). An arclength parametrized curve $\gamma(s)=\Y(u(s),v(s))$, with $\kappa_\gamma>0$, is a geodesic of $M$ if and only if there is a differentiable function $\omega$ such that
\begin{align}
u'(s)\,(v(s)-u(s))\,\kappa_\alpha(u(s)) &= \sin\omega(s),\label{p2.1}\\
v'(s) &= \cos\omega(s),\label{p2.2}\\
u'(s)\,\kappa_\alpha(u(s)) &= -\omega'(s).\label{p2.3}
\end{align}
Moreover, the curvature and torsion of $\gamma$ are given by (\ref{kg2}) and (\ref{tg2}), respectively.
\end{proposition}

\begin{align}
u'(s)\,(v(s)-u(s))\,\kappa_\alpha(u(s)) &= \sin\omega(s),\\
v'(s) &= \cos\omega(s),\\
u'(s)\,\kappa_\alpha(u(s)) &= -\omega'(s).
\end{align}

\begin{theorem}\label{th1}
Let $\gamma(s)$ be an arclength parametrized curve fully immersed in $\R3$. If $\gamma$ is a proper concircular helix, then there exists a rectifying curve $\alpha$ such that $\gamma$ is (congruent to) a geodesic of the tangent surface to $\alpha$.
\end{theorem}

\begin{proof}
Let us consider the following functions:
\begin{align}
\omega(s)= & -\arccot\rho(s), \label{w teo}\\
u(s)= & -\frac{1}{m}\left(\kappa_\gamma(s)\,\frac{(1+\rho(s)^2)^{3/2}}{\rho'(s)}+n\right), \label{u teo}\\
v(s)= & u(s)-\frac{\sin\omega(s)}{\omega'(s)}, \label{v teo}
\end{align}
for a constant $n$. Here, $m$ and $\rho(s)$ are given in Theorem \ref{teoconc1}.

Let $\alpha$ be the rectifying curve determined by the curvature function
\begin{equation}\label{teo kalpha}
  \kappa_\alpha(t)=-\frac{\omega'(u^{-1}(t))}{u'(u^{-1}(t))}
\end{equation}
and whose torsion satisfies $(\tau_\alpha/\kappa_\alpha)(u)=m\,u+n$.

On the tangent surface to $\alpha$, let us take the curve $\tgamma(s)=\Y(u(s),v(s))$. From (\ref{teo kalpha}) we get (\ref{p2.3}), and then equation (\ref{v teo}) leads to (\ref{p2.1}).

By derivating (\ref{u teo}) and using (\ref{conc4}) we obtain
\[
u'(s)=\sqrt{1+\rho(s)^2}\,\frac{\rho''(s)}{\rho'(s)^2},
\]
and then we get
\[
v'(s)=\frac{\rho''(s)\sqrt{1+\rho(s)^2}}{\rho'(s)^2}+ \left(\frac{\sqrt{1+\rho(s)^2}}{\rho'(s)}\right)'
   =\frac{\rho(s)}{\sqrt{1+\rho^2(s)}}=\cos\omega(s).
\]
In conclusion, by Proposition \ref{p3} we get $\tgamma$ is a geodesic in the tangent surface to a rectifying curve, so it is a concircular helix.

On the other hand, we have
\begin{equation*}
  \frac{\tau_\alpha}{\kappa_\alpha}(u(s))=m\,u(s)+n=-\kappa_\gamma(s)\,\frac{(1+\rho(s)^2)^{3/2}}{\rho'(s)}.
\end{equation*}
Since the curvature and torsion of $\tilde\gamma$ satisfy (\ref{kg2}) y (\ref{tg2}), respectively, we get
\begin{align*}
  \kappa_{\tgamma}(s)= &  \frac{-1}{\sqrt{1+\rho(s)^2}}\,\frac{\rho'(s)}{1+\rho(s)^2}\,\frac{-(1+\rho(s)^2)^{3/2}}{\rho'(s)}\,\kappa_\gamma(s)=\kappa_\gamma(s),\\
  \tau_{\tgamma}(s)= &  -\frac{\rho(s)}{\sqrt{1+\rho(s)^2}}\,\frac{\rho'(s)}{1+\rho(s)^2}\,\frac{-(1+\rho(s)^2)^{3/2}}{\rho'(s)}\,\kappa_\gamma(s)=\tau_\gamma(s),
\end{align*}
which shows that $\gamma$ and $\tgamma$ are congruent curves.
\end{proof}

\end{document}
